% !PRINTED 248
\weberpage{248}
%<<CONT
real, but positive when $\alpha$ and $\gamma$ are conjugate (and then, because $\alpha+\gamma=0$, purely imaginary). From this it follows:

The biquadratic equation has two pairs of equal real roots when
\weq{8}{D = 0, \quad a < 0, \quad a^2-4c = 0}
and two pairs of equal imaginary roots when
\weq{9}{D = 0, \quad a > 0, \quad a^2-4c = 0.}

But if $\alpha$ is different from $\beta$, then $\gamma$ must be real, and (7) shows that when $\alpha$ and $\beta$ are real, $a$ is negative and $a^2-4c$ is positive; that, on the other hand, when $\alpha$ and $\beta$ are conjugate imaginary, either $a$ must be positive or $a^2-4c$ negative, since in this case $(\alpha-\beta)^2$ is negative, and if $2(\alpha+\beta)^2+(\alpha-\beta)^2$ is positive, then certainly $8(\alpha+\beta)^2+(\alpha-\beta)^2$ must also be positive.

We can therefore supplement (6) by saying:

\gesperrt{The necessary and sufficient condition for the existence of four real roots, of which two (but not more) may also be equal to one another, is}
\weq{10}{D \geqq 0, \quad a < 0, \quad a^2-4c > 0.}

\gesperrt{The necessary and sufficient condition for three equal roots,} which are necessarily all real, is by (2)
\[
u^2 = v^2 = w^2,
\]
hence the cubic resolvent (3) must have three equal roots, and for this, by \S.\,77 (1), the necessary and sufficient conditions are
\weq{11}{\begin{aligned} a^2+12c &= 0\\ 2a^3-8ac+9b^2 &= 0. \end{aligned}}
By \S.\,47, (15), (16) the two invariants $A$ and $B$ of the biquadratic equation are
\[
A = a^2+12c, \quad B = 2a^3-72ac+27b^2.
\]
Hence we can also write the equations (11) thus
\[
A = 0, \quad B+4aA = 0,
\]
so that (11) is equivalent to
\[
A = 0, \quad B = 0.
\]
By elimination of $c$ one can also derive from (11) the equation:
\weq{12}{8a^3+27b^2 = 0,}
